\par\medskip\noindent\textbf{Cómo leer la salida.}
{\small\setlength{\tabcolsep}{4pt}\renewcommand{\arraystretch}{1.15}
\begin{xltabular}{\linewidth}{@{}>{\raggedright\arraybackslash}p{0.24\linewidth}>{\raggedleft\arraybackslash}p{0.15\linewidth}>{\raggedright\arraybackslash}X@{}}
\toprule Indicador & Valor & Interpretación \\ \midrule \endfirsthead
\toprule Indicador & Valor & Interpretación \\ \midrule \endhead
\bottomrule \endlastfoot
Meses en recesión & 199 de 1109 & Meses en recesión según el NBER. \\
Pr media en recesión / expansión & 0,437 / 0,079 & El régimen volátil es mucho más probable en recesión. \\
Tabla de clasificación & 852; 58; 110; 89 & Coincidencias y discrepancias entre el régimen estimado y el NBER. \\
Pr(volátil $\mid$ recesión) & 0,447 & Menos de la mitad de los meses de recesión son volátiles: no toda recesión agita la bolsa. \\
Pr(volátil $\mid$ expansión) & 0,064 & Pocas falsas alarmas. \\
Correlación & 0,466 & Relación moderada entre la probabilidad suavizada y el indicador de recesión. \\
Episodios: inicio, fin, meses & p. ej. 1929-05 a 1934-07, 63 & Períodos volátiles de al menos 6 meses; 2000-2001 y 2008-2009 aparecen claramente. \\
\texttt{rend. medio}, \texttt{desv. est.} & $-4{,}62$ a 0,13; 5,3 a 12,8 & Rendimiento y volatilidad en cada episodio. \\
Meses en recesión del episodio & p. ej. 43 de 63 & Casi todos los episodios se solapan con recesiones. \\
Pr media por recesión & 0,96 (1929) a 0,002 (1953) & Las recesiones de 1945--1961 y 1980--1991 no generaron volatilidad bursátil; las de 1929, 1937, 1973, 2001 y 2008 sí. \\
\end{xltabular}}
